\documentclass[11pt,a4paper,parskip=half,headsepline=0.6pt,footsepline=0.4pt]{scrartcl}

% ---------------------------------------------------------------
% Packages & Core Configuration
% ---------------------------------------------------------------
\usepackage[utf8]{inputenc}
\usepackage[T1]{fontenc}
\usepackage[scaled=0.92]{helvet}
\renewcommand{\familydefault}{\sfdefault}
\usepackage[scaled=0.88]{beramono}

\usepackage{amsmath,amssymb,amsfonts}
\usepackage{graphicx}
\usepackage[table,dvipsnames]{xcolor}
\usepackage{booktabs,tabularx,array,multirow,colortbl}
\usepackage[margin=15mm,top=18mm,bottom=18mm,headheight=15pt]{geometry}
\usepackage{scrlayer-scrpage}
\usepackage{lastpage}
\usepackage{siunitx}
\usepackage{fontawesome5}
\usepackage{enumitem}
\usepackage{microtype}

\usepackage{tikz}
\usetikzlibrary{positioning,calc,shapes.geometric,patterns,decorations.pathmorphing,arrows.meta,backgrounds}

\usepackage[many]{tcolorbox}

% ---------------------------------------------------------------
% Custom Color Palette (Modern 3-Color Scientific Theme)
% ---------------------------------------------------------------
\definecolor{navyprimary}{RGB}{24,43,73}
\definecolor{tealaccent}{RGB}{0,136,145}
\definecolor{coralhighlight}{RGB}{224,86,36}
\definecolor{darkgraytext}{RGB}{40,44,52}
\definecolor{lightbg}{RGB}{245,247,250}
\definecolor{tableheader}{RGB}{230,236,245}
\definecolor{bordergray}{RGB}{205,212,223}
\definecolor{cardbg}{RGB}{252,253,255}

% ---------------------------------------------------------------
% Header & Footer Settings (scrlayer-scrpage)
% ---------------------------------------------------------------
\pagestyle{scrheadings}
\ihead{} \chead{} \ohead{} \ifoot{} \cfoot{} \ofoot{}
\ihead{\footnotesize \textcolor{navyprimary}{\textbf{\faFlask\ Advanced Optics \& Electronics Lab} \ \textbar\ Exp.\ \#04}}
\chead{}
\ohead{\footnotesize \textcolor{darkgraytext}{Station: \textbf{B-07} \ \textbar\ Vertex Institute of Tech.}}
\ifoot{\footnotesize \textcolor{darkgraytext}{\faAtom\ PHYS-402: Quantum Photonics \& Circuit Metrology}}
\cfoot{}
\ofoot{\footnotesize \textcolor{darkgraytext}{Page \textbf{\thepage} of \textbf{\pageref{LastPage}}}}

\addtokomafont{headsepline}{\color{tealaccent}}
\addtokomafont{footsepline}{\color{bordergray}}

% ---------------------------------------------------------------
% Custom tcolorbox Styles
% ---------------------------------------------------------------
\tcbset{
  labbox/.style={
    enhanced,
    colback=lightbg,
    colframe=navyprimary,
    fonttitle=\bfseries\sffamily,
    coltitle=white,
    boxrule=0.8pt,
    arc=3pt,
    top=6pt,bottom=6pt,left=8pt,right=8pt,
    shadow={2pt}{-2pt}{0pt}{black!15}
  },
  calloutbox/.style={
    enhanced,
    colback=coralhighlight!5,
    colframe=coralhighlight,
    fonttitle=\bfseries\sffamily,
    coltitle=coralhighlight!80!black,
    top=4pt,bottom=4pt,left=8pt,right=8pt,
    boxrule=0.8pt,
    arc=2pt,
    leftrule=3.5pt
  },
  accentcard/.style={
    enhanced,
    colback=cardbg,
    colframe=tealaccent,
    boxrule=0.6pt,
    arc=3pt,
    top=4pt,bottom=4pt,left=7pt,right=7pt
  }
}

% ---------------------------------------------------------------
% Section Title Formatting
% ---------------------------------------------------------------
\addtokomafont{section}{\color{navyprimary}\Large\bfseries}
\addtokomafont{subsection}{\color{tealaccent}\large\bfseries}
\addtokomafont{subsubsection}{\color{darkgraytext}\normalfont\bfseries}

\renewcommand{\thesection}{\Roman{section}}
\renewcommand{\thesubsection}{\thesection.\arabic{subsection}}

% Tighten spacing around sections for clean 3-page layout
\RedeclareSectionCommand[
  beforeskip=-0.6ex plus -0.2ex minus -0.1ex,
  afterskip=0.3ex plus 0.1ex
]{section}
\RedeclareSectionCommand[
  beforeskip=-0.5ex plus -0.2ex minus -0.1ex,
  afterskip=0.2ex plus 0.1ex
]{subsection}

% ---------------------------------------------------------------
% Document Content
% ---------------------------------------------------------------
\begin{document}

% ===============================================================
% DOCUMENT HEADER / TITLE BLOCK
% ===============================================================
\begin{tcolorbox}[labbox, colback=navyprimary, colframe=navyprimary, top=7pt, bottom=7pt, left=8pt, right=8pt]
\begin{minipage}[c]{0.68\linewidth}
  \textcolor{white}{\small \textbf{VERTEX INSTITUTE OF TECHNOLOGY} \ \textbar\ Department of Applied Physics}\\
  \vspace{1pt}
  {\Large \bfseries \textcolor{white}{Laser Diode Dynamics, Interferometry \& TIA Circuit Record}} \\
  \vspace{2pt}
  \textcolor{tealaccent!30}{\small \faMicrochip\ Course Code: \textbf{PHYS-402} \ \textbar\ Module: Electro-Optic Measurement Systems}
\end{minipage}%
\hfill
\begin{minipage}[c]{0.29\linewidth}
  \begin{tcolorbox}[enhanced, colback=white, colframe=tealaccent, boxrule=0.8pt, arc=3pt, top=3pt, bottom=3pt, left=4pt, right=4pt]
    \centering
    \textcolor{navyprimary}{\footnotesize \textbf{EXPERIMENT RECORD}}\\
    \vspace{1pt}
    \textcolor{coralhighlight}{\large \textbf{\# EXP-2026-04}}\\
    \vspace{1pt}
    \textcolor{darkgraytext}{\tiny Status: \textbf{VERIFIED \& PASSED}}
  \end{tcolorbox}
\end{minipage}
\end{tcolorbox}

\vspace{-4pt}

% Metadata Bar
\begin{tcolorbox}[accentcard, colback=lightbg, colframe=bordergray, top=3pt, bottom=3pt, left=6pt, right=6pt]
\small
\begin{tabularx}{\linewidth}{>{\raggedright\arraybackslash}X >{\raggedright\arraybackslash}X >{\raggedright\arraybackslash}X >{\raggedright\arraybackslash}X}
  \textbf{Lead Exp.:} Elena Rostova & \textbf{Date:} 2026-07-24 & \textbf{Station:} Bench B-07 & \textbf{Ambient Temp:} $21.4\,^{\circ}\text{C}$ \\
  \textbf{Lab Partner:} Julian Thorne & \textbf{Supervisor:} Dr. Vance & \textbf{Laser Class:} Class 3B (650nm) & \textbf{Rel. Humidity:} $44\,\%$
\end{tabularx}
\end{tcolorbox}

\vspace{1pt}

% ===============================================================
% SECTION I: EXECUTIVE SUMMARY & OBJECTIVES
% ===============================================================
\section{Executive Summary \& Experimental Objectives}

This laboratory investigation presents a complete electro-optical characterization of a Distributed Feedback (DFB) Semiconductor Laser Diode ($\lambda_0 = 650\,\text{nm}$), an ultra-sensitive Michelson Interferometric Displacement System, and a wideband Photodiode Transimpedance Amplifier (TIA) signal conditioning stage.

\vspace{1pt}

\begin{minipage}[t]{0.48\linewidth}
  \begin{tcolorbox}[accentcard, title={\small \textcolor{navyprimary}{\faBullseye\ Key Experimental Objectives}}, colframe=navyprimary!40]
    \small
    \begin{enumerate}[leftmargin=1.2em, itemsep=1pt, topsep=1pt]
      \item Measure $L-I-V$ characteristic curves to determine threshold current $I_{th}$ and slope efficiency $\eta_d$.
      \item Quantify fringe visibility $V$ as a function of optical path difference $\Delta L$ in a Michelson setup.
      \item Optimize a high-speed TIA circuit for gain-bandwidth ($f_{-3dB} \ge 12\,\text{MHz}$) with phase margin $\phi_m \ge 60^\circ$.
    \end{enumerate}
  \end{tcolorbox}
\end{minipage}%
\hfill
\begin{minipage}[t]{0.48\linewidth}
  \begin{tcolorbox}[accentcard, title={\small \textcolor{navyprimary}{\faSlidersH\ Calibrated Test Instruments}}, colframe=navyprimary!40]
    \small
    \begin{tabularx}{\linewidth}{>{\raggedright\arraybackslash}X l c c}
      \toprule
      \textbf{Device} & \textbf{Model} & \textbf{S/N} & \textbf{Uncert.} \\
      \midrule
      Power Meter & ThorLabs & OP-9842 & $\pm 0.8\%$ \\
      Oscilloscope & Keysight & MY5412 & $\pm 0.5\%$ \\
      Piezo Ctrl & ThorPZT & PZ-4410 & $\pm 1.2\,\text{nm}$ \\
      SourceMeter & Keithley & SM-8821 & $\pm 0.05\%$ \\
      \bottomrule
    \end{tabularx}
  \end{tcolorbox}
\end{minipage}

\par \vspace{2pt}

% Safety Callout
\begin{tcolorbox}[calloutbox]
  \textcolor{coralhighlight}{\textbf{\faExclamationTriangle\ Safety Precautions \& ESD Operating Rules:}} \\
  \footnotesize
  1.~Class 3B laser radiation hazard ($P_{max} = 15\,\text{mW}$ at $650\,\text{nm}$). Alignment must be performed strictly using certified OD4+ laser safety eyewear. 2.~PIN photodiode and OPA847 op-amp inputs are extremely sensitive to electrostatic discharge ($V_{ESD, max} = 250\,\text{V}$). Wrist grounding straps were worn throughout all circuit assemblies.
\end{tcolorbox}

\vspace{1pt}

% ===============================================================
% SECTION II: THEORETICAL FOUNDATIONS
% ===============================================================
\section{Theoretical Foundations \& Governing Equations}

The optical output power $P_{opt}$ of the semiconductor laser diode above threshold current $I_{th}$ is governed by the differential quantum efficiency $\eta_d$:
\begin{equation}
  P_{opt}(I) = \eta_d \frac{\hbar \omega}{q} \left( I - I_{th} \right) \quad \text{for } I > I_{th}
\end{equation}
where $\hbar \omega = \frac{hc}{\lambda_0}$ represents the photon energy and $q = 1.602 \times 10^{-19}\,\text{C}$ is the elementary electronic charge.

In the Michelson interferometer configuration, the superposition of the reference beam $E_1$ and signal beam $E_2$ yields a total detected optical intensity $I_{det}(\delta)$:
\begin{equation}
  I_{det}(\delta) = I_1 + I_2 + 2\sqrt{I_1 I_2} \, \gamma(\tau) \cos\left( \frac{4\pi \Delta d}{\lambda_0} + \phi_0 \right)
\end{equation}
where $\Delta d$ is the mirror displacement, $\gamma(\tau)$ is the complex degree of temporal coherence, and the interference fringe visibility $V$ is defined as:
\begin{equation}
  V = \frac{I_{max} - I_{min}}{I_{max} + I_{min}} = \frac{2\sqrt{I_1 I_2}}{I_1 + I_2} \, \big|\gamma(\tau)\big|
\end{equation}

For the front-end optical detection circuit, the photocurrent $I_{pd} = \mathcal{R} \cdot P_{opt}$ (where $\mathcal{R} = 0.45\,\text{A/W}$ is photodiode responsivity) is converted to output voltage $V_{out}$ via a transimpedance amplifier. The $-3\,\text{dB}$ bandwidth $f_{-3dB}$ and optimal feedback capacitance $C_f$ to prevent peaking are given by Butterworth filter criteria:
\begin{equation}
  C_f = \sqrt{\frac{C_{in}}{2\pi R_f f_T}}, \qquad f_{-3dB} = \sqrt{\frac{f_T}{2\pi R_f C_{in}}}
\end{equation}
where $C_{in} = C_{pd} + C_{diff} + C_{cm}$ is the total input capacitance and $f_T$ is the Gain-Bandwidth Product (GBWP) of the operational amplifier.

\newpage

% ===============================================================
% SECTION III: EXPERIMENTAL SETUPS & SCHEMATICS
% ===============================================================
\section{Experimental Setups \& Schematics}

\begin{minipage}[t]{0.48\linewidth}
  \centering
  \textbf{\small Figure 1: Michelson Interferometer Setup} \\ \vspace{2pt}
  \begin{tikzpicture}[scale=0.68, every node/.style={transform shape}]
    % Background grid
    \draw[lightbg, fill=lightbg, rounded corners=4pt] (-0.3,-0.3) rectangle (6.2,5.2);
    
    % Laser Source
    \draw[fill=navyprimary!90, draw=navyprimary, text=white, font=\bfseries\tiny] (0,2) rectangle (1.2,2.8) node[pos=0.5]{650nm Laser};
    
    % Beam Splitter
    \draw[fill=tealaccent!20, draw=tealaccent, thick] (2.8,1.9) -- (3.7,2.8) -- (3.5,3.0) -- (2.6,2.1) -- cycle;
    \node[font=\tiny\bfseries, color=navyprimary] at (3.8,1.6) {50:50 BS};
    
    % Fixed Mirror (Top)
    \draw[fill=darkgraytext!80, draw=black, thick] (2.7,4.5) rectangle (3.7,4.7);
    \draw[pattern=north east lines, pattern color=gray] (2.7,4.7) rectangle (3.7,4.9);
    \node[font=\tiny\bfseries] at (3.2,5.05) {Fixed Mirror $M_1$};
    
    % Piezo Mirror (Right)
    \draw[fill=coralhighlight!70, draw=coralhighlight!90, thick] (5.4,1.9) rectangle (5.6,2.9);
    \draw[fill=gray!40] (5.6,2.0) rectangle (5.8,2.8);
    \node[font=\tiny\bfseries, color=coralhighlight!90!black, align=center] at (5.6,1.4) {Piezo $M_2$\\ ($\Delta d$)};
    
    % Photodetector (Bottom)
    \draw[fill=tealaccent!80, draw=tealaccent!90, text=white, font=\tiny\bfseries] (2.7,0.1) rectangle (3.7,0.7) node[pos=0.5]{Photodiode};
    
    % Laser Beams
    % Main Beam from Laser to BS
    \draw[red!90!black, ultra thick, ->, >=stealth] (1.2,2.4) -- (3.05,2.4);
    % Transmitted Beam to M2
    \draw[red!90!black, ultra thick, ->, >=stealth] (3.05,2.4) -- (5.4,2.4);
    \draw[red!70!black, thick, dashed, <-, >=stealth] (3.15,2.5) -- (5.4,2.5);
    % Reflected Beam to M1
    \draw[red!90!black, ultra thick, ->, >=stealth] (3.15,2.45) -- (3.15,4.5);
    \draw[red!70!black, thick, dashed, <-, >=stealth] (3.05,2.45) -- (3.05,4.5);
    % Combined Beam to Photodetector
    \draw[red!90!black, double, ultra thick, ->, >=stealth] (3.1,2.35) -- (3.1,0.7);
    
    % Labels
    \node[font=\tiny, color=red!80!black] at (2.1,2.6) {$E_{in}$};
    \node[font=\tiny, color=red!80!black] at (4.3,2.7) {$E_2$};
    \node[font=\tiny, color=red!80!black] at (3.5,3.6) {$E_1$};
  \end{tikzpicture}
\end{minipage}%
\hfill
\begin{minipage}[t]{0.48\linewidth}
  \centering
  \textbf{\small Figure 2: Photodiode TIA Signal Conditioning Circuit} \\ \vspace{2pt}
  \begin{tikzpicture}[scale=0.68, every node/.style={transform shape}]
    % Background card
    \draw[lightbg, fill=lightbg, rounded corners=4pt] (-0.5,-0.3) rectangle (6.2,5.2);
    
    % Photodiode Symbol
    \draw[thick] (0,2) -- (0.6,2);
    \draw[thick] (0.6,1.6) -- (0.6,2.4);
    \draw[fill=navyprimary!20, thick] (0.6,2) -- (1.1,2.4) -- (1.1,1.6) -- cycle;
    \draw[thick] (1.1,1.6) -- (1.1,2.4);
    \draw[thick] (1.1,2) -- (1.8,2);
    % Light arrows onto Photodiode
    \draw[coralhighlight, thick, ->, >=stealth] (0.3,2.7) -- (0.7,2.2);
    \draw[coralhighlight, thick, ->, >=stealth] (0.5,2.9) -- (0.9,2.4);
    \node[font=\tiny\bfseries, color=coralhighlight!90!black] at (0.3,3.1) {$P_{opt}$};
    \node[font=\tiny] at (0.8,1.3) {PIN PD};
    
    % Ground at Anode
    \draw[thick] (0,2) -- (0,1.2);
    \draw[thick] (-0.2,1.2) -- (0.2,1.2);
    \draw[thick] (-0.14,1.05) -- (0.14,1.05);
    \draw[thick] (-0.07,0.9) -- (0.07,0.9);
    
    % Op Amp Triangle
    \draw[thick, fill=white] (2.2,1.3) -- (3.8,2.1) -- (2.2,2.9) -- cycle;
    \node[font=\tiny\bfseries] at (2.4,2.6) {$-$};
    \node[font=\tiny\bfseries] at (2.4,1.6) {$+$};
    \node[font=\scriptsize\bfseries, color=navyprimary] at (2.9,2.1) {OPA847};
    
    % Non-inverting input ground
    \draw[thick] (1.8,1.6) -- (2.2,1.6);
    \draw[thick] (1.8,1.6) -- (1.8,1.0);
    \draw[thick] (1.6,1.0) -- (2.0,1.0);
    \draw[thick] (1.7,0.85) -- (1.9,0.85);
    
    % Inverting input connection
    \draw[thick] (1.8,2.0) -- (1.8,2.6) -- (2.2,2.6);
    
    % Feedback loop (Rf and Cf in parallel)
    \draw[thick] (1.8,2.6) -- (1.8,4.2) -- (2.3,4.2);
    \draw[thick] (3.3,4.2) -- (4.4,4.2) -- (4.4,2.1) -- (3.8,2.1);
    
    % Feedback Resistor Rf (Top branch)
    \draw[thick] (2.3,4.5) -- (3.3,4.5);
    \draw[fill=tealaccent!15, draw=tealaccent, thick] (2.5,4.3) rectangle (3.1,4.7);
    \node[font=\tiny\bfseries, color=navyprimary] at (2.8,4.9) {$R_f = 100\,\text{k}\Omega$};
    
    % Feedback Capacitor Cf (Bottom branch)
    \draw[thick] (2.3,3.9) -- (2.7,3.9);
    \draw[thick] (2.7,3.7) -- (2.7,4.1);
    \draw[thick] (2.9,3.7) -- (2.9,4.1);
    \draw[thick] (2.9,3.9) -- (3.3,3.9);
    \node[font=\tiny\bfseries, color=navyprimary] at (2.8,3.4) {$C_f = 1.5\,\text{pF}$};
    
    \draw[thick] (1.8,4.2) -- (1.8,4.5);
    \draw[thick] (1.8,4.2) -- (1.8,3.9);
    \draw[thick] (3.3,4.2) -- (3.3,4.5);
    \draw[thick] (3.3,4.2) -- (3.3,3.9);
    
    % Output terminal
    \draw[thick, ->, >=stealth] (3.8,2.1) -- (5.2,2.1);
    \node[circle, fill=navyprimary, inner sep=1pt] at (5.2,2.1) {};
    \node[font=\small\bfseries, color=navyprimary] at (5.7,2.1) {$V_{out}$};
  \end{tikzpicture}
\end{minipage}

\vspace{1pt}

% ===============================================================
% SECTION IV: EXPERIMENTAL DATA LOGS
% ===============================================================
\section{Tabulated Experimental Measurement Logs}

\subsection{Laser Diode Light--Current--Voltage ($L-I-V$) Characteristic Sweep}

Laser Junction Temperature maintained at $T_j = (25.0 \pm 0.1)^\circ\text{C}$ via TEC controller. Optical power measured with ThorLabs PM100D at distance $d = 15.0\,\text{cm}$.

\vspace{1pt}

\begin{table}[h!]
\centering
\footnotesize
\caption{Measured Laser Diode $L-I-V$ Experimental Data}
\label{tab:liv_data}
\vspace{1pt}
\begin{tabularx}{\linewidth}{>{\centering\arraybackslash}p{2.2cm} >{\centering\arraybackslash}X >{\centering\arraybackslash}X >{\centering\arraybackslash}X >{\centering\arraybackslash}X >{\centering\arraybackslash}p{3.2cm}}
  \toprule
  \rowcolor{tableheader}
  \textbf{Run No.} & \textbf{Drive Current $I_d$ (mA)} & \textbf{Forward Voltage $V_d$ (V)} & \textbf{Optical Power $P_{opt}$ (mW)} & \textbf{Responsivity $\mathcal{R}_{eff}$ (A/W)} & \textbf{Operating State} \\
  \midrule
  01 & 2.00 & 1.240 & 0.002 & 0.001 & LED Spontaneous \\
  02 & 6.00 & 1.512 & 0.011 & 0.002 & LED Spontaneous \\
  03 & 10.00 & 1.710 & 0.028 & 0.003 & LED Spontaneous \\
  04 & 14.00 & 1.845 & 0.052 & 0.004 & Amplified Spont. \\
  05 & 17.50 & 1.928 & 0.095 & 0.005 & Near Threshold \\
  \rowcolor{tealaccent!10}
  06 (Threshold) & \textbf{18.65} & \textbf{1.954} & \textbf{0.210} & \textbf{0.011} & \textbf{Lasing Onset ($I_{th}$)} \\
  07 & 22.00 & 2.021 & 1.840 & 0.084 & Stimulated Emission \\
  08 & 26.00 & 2.095 & 3.910 & 0.150 & Stimulated Emission \\
  09 & 30.00 & 2.162 & 5.980 & 0.199 & Linear Lasing \\
  10 & 34.00 & 2.224 & 8.040 & 0.236 & Linear Lasing \\
  11 & 38.00 & 2.281 & 10.120 & 0.266 & Linear Lasing \\
  12 & 42.00 & 2.335 & 12.180 & 0.290 & High Power Lasing \\
  \bottomrule
\end{tabularx}
\end{table}

\vspace{-6pt}

\subsection{Michelson Interferometer Displacement vs. Fringe Visibility}

Fringe intensity measurements recorded by scanning Piezo Actuator $M_2$ with step size $\Delta V_{PZT} = 1.0\,\text{V}$ (displacement $\approx 85.2\,\text{nm/V}$).

\vspace{1pt}

\begin{table}[h!]
\centering
\footnotesize
\caption{Interferometric Fringe Intensity \& Calculated Visibility}
\label{tab:visibility_data}
\vspace{1pt}
\begin{tabularx}{\linewidth}{>{\centering\arraybackslash}c >{\centering\arraybackslash}X >{\centering\arraybackslash}X >{\centering\arraybackslash}X >{\centering\arraybackslash}X >{\centering\arraybackslash}c}
  \toprule
  \rowcolor{tableheader}
  \textbf{Arm Difference $\Delta L$ (mm)} & \textbf{$I_{max}$ Voltage (mV)} & \textbf{$I_{min}$ Voltage (mV)} & \textbf{Peak-to-Peak $\Delta V$ (mV)} & \textbf{Visibility $V$} & \textbf{Coherence Status} \\
  \midrule
  0.0 (Equal Path) & $845.0 \pm 2.1$ & $12.5 \pm 0.8$ & 832.5 & \textbf{0.971} & Peak Coherence \\
  5.0 & $832.0 \pm 2.0$ & $24.0 \pm 0.9$ & 808.0 & \textbf{0.944} & High Coherence \\
  15.0 & $795.0 \pm 2.5$ & $58.0 \pm 1.2$ & 737.0 & \textbf{0.864} & Good Visibility \\
  30.0 & $720.0 \pm 3.0$ & $125.0 \pm 1.8$ & 595.0 & \textbf{0.704} & Partial Coherence \\
  50.0 & $610.0 \pm 3.2$ & $215.0 \pm 2.1$ & 395.0 & \textbf{0.479} & Moderate Decay \\
  75.0 & $480.0 \pm 3.8$ & $310.0 \pm 2.5$ & 170.0 & \textbf{0.215} & Near Coherence Length \\
  100.0 & $415.0 \pm 4.0$ & $370.0 \pm 3.1$ & 45.0 & \textbf{0.057} & Spatial Washout \\
  \bottomrule
\end{tabularx}
\end{table}

\newpage

% ===============================================================
% SECTION V: UNCERTAINTY ANALYSIS & LINEAR FIT
% ===============================================================
\section{Uncertainty Analysis \& Statistical Evaluation}

\subsection{Error Propagation Analysis}

Uncertainties in optical slope efficiency $\eta_d = \frac{q}{\hbar \omega} \frac{\Delta P_{opt}}{\Delta I}$ were evaluated using standard Gaussian error propagation:
\begin{equation}
  \sigma_{\eta_d} = \eta_d \sqrt{ \left( \frac{\sigma_{\Delta P}}{\Delta P} \right)^2 + \left( \frac{\sigma_{\Delta I}}{\Delta I} \right)^2 + \left( \frac{\sigma_{\lambda}}{\lambda_0} \right)^2 }
\end{equation}

\begin{table}[h!]
\centering
\small
\caption{Uncertainty Budget and Component Breakdown}
\label{tab:uncertainty_budget}
\vspace{1pt}
\begin{tabularx}{\linewidth}{X >{\centering\arraybackslash}p{2.5cm} >{\centering\arraybackslash}p{3.0cm} >{\centering\arraybackslash}p{2.8cm} >{\centering\arraybackslash}p{2.5cm}}
  \toprule
  \rowcolor{tableheader}
  \textbf{Parameter / Quantity} & \textbf{Nominal Value} & \textbf{Absolute Error ($\sigma_x$)} & \textbf{Relative Error (\%)} & \textbf{Primary Source} \\
  \midrule
  Drive Current $I_d$ & $30.00\,\text{mA}$ & $\pm 0.015\,\text{mA}$ & $0.05\%$ & Keithley SMU \\
  Laser Wavelength $\lambda_0$ & $650.2\,\text{nm}$ & $\pm 0.4\,\text{nm}$ & $0.06\%$ & Spectrometer \\
  Optical Power $P_{opt}$ & $5.980\,\text{mW}$ & $\pm 0.048\,\text{mW}$ & $0.80\%$ & ThorLabs Meter \\
  Piezo Position $\Delta d$ & $426.0\,\text{nm}$ & $\pm 5.1\,\text{nm}$ & $1.20\%$ & PZT Hysteresis \\
  Calculated Efficiency $\eta_d$ & $0.518\,\text{W/A}$ & $\pm 0.006\,\text{W/A}$ & $1.16\%$ & Combined Prop. \\
  \bottomrule
\end{tabularx}
\end{table}

\vspace{1pt}

\subsection{Linear Regression Fit Results (Laser L-I Lasing Region)}

Linear regression analysis performed on Data Points 07 to 12 ($I_d = 22.0\,\text{mA}$ to $42.0\,\text{mA}$):
\begin{align}
  P_{opt}(I_d) &= S \cdot I_d + P_0 \\[2pt]
  \text{Fitted Slope } S &= (0.517 \pm 0.003)\,\text{mW/mA} = (0.517 \pm 0.003)\,\text{W/A} \\
  \text{Calculated Threshold } I_{th} &= (18.58 \pm 0.12)\,\text{mA} \\
  \text{Coefficient of Determination } R^2 &= 0.9994 \quad (\text{Excellent linearity})
\end{align}

\vspace{1pt}

% Summary Performance Cards
\begin{minipage}[t]{0.48\linewidth}
  \begin{tcolorbox}[accentcard, title={\small \textcolor{navyprimary}{\faCheckCircle\ Measured Key Parameters}}, colframe=tealaccent]
    \small
    \begin{itemize}[leftmargin=1.2em, itemsep=1.5pt]
      \item \textbf{Threshold Current $I_{th}$:} $(18.58 \pm 0.12)\,\text{mA}$
      \item \textbf{Slope Efficiency $\eta_d$:} $(0.517 \pm 0.003)\,\text{W/A}$
      \item \textbf{Coherence Length $L_c$:} $(68.4 \pm 2.2)\,\text{mm}$
      \item \textbf{TIA $-3\text{dB}$ Bandwidth:} $14.2\,\text{MHz}$
    \end{itemize}
  \end{tcolorbox}
\end{minipage}%
\hfill
\begin{minipage}[t]{0.48\linewidth}
  \begin{tcolorbox}[accentcard, title={\small \textcolor{navyprimary}{\faChartLine\ Circuit Performance Summary}}, colframe=tealaccent]
    \small
    \begin{itemize}[leftmargin=1.2em, itemsep=1.5pt]
      \item \textbf{Transimpedance Gain:} $100\,\text{kV/A}$ ($100\,\text{k}\Omega$)
      \item \textbf{Input Voltage Noise:} $0.98\,\text{nV}/\sqrt{\text{Hz}}$
      \item \textbf{Phase Margin $\phi_m$:} $64.2^\circ$ (Butterworth response)
      \item \textbf{Total Output Offset:} $< 1.2\,\text{mV}$
    \end{itemize}
  \end{tcolorbox}
\end{minipage}

\vspace{4pt}

% ===============================================================
% SECTION VI: DISCUSSION & VERIFICATION
% ===============================================================
\section{Discussion, Observations \& Lab Verification}

\subsection{Physical Insights \& Experimental Deviations}
1. \textbf{Threshold Behavior}: The experimental threshold current of $18.58\,\text{mA}$ matches the manufacturer datasheet rating ($18.50\,\text{mA}$) within $0.4\%$. A minor thermal drift was observed when operating above $35\,\text{mA}$, which was suppressed by adjusting the TEC loop PID parameters.
2. \textbf{Fringe Visibility Decay}: The fringe visibility $V(\Delta L)$ decayed from $0.971$ at zero path difference down to $0.057$ at $\Delta L = 100\,\text{mm}$. The temporal coherence length $L_c \approx 68.4\,\text{mm}$ yields a spectral linewidth estimate $\Delta \lambda = \lambda_0^2 / L_c \approx 6.18\,\text{pm}$, consistent with single-mode DFB laser characteristics.
3. \textbf{TIA Compensation}: Without the feedback capacitor $C_f = 1.5\,\text{pF}$, the circuit displayed $3.2\,\text{dB}$ of peaking at $16\,\text{MHz}$ due to photodiode capacitance $C_{pd} = 12\,\text{pF}$. Installing $C_f$ successfully restored a smooth maximally-flat frequency response.

\vspace{4pt}

% Formal Verification Sign-off Box
\begin{tcolorbox}[labbox, colback=white, colframe=navyprimary, title={\small \faSignature\ Formal Laboratory Verification \& Approval Register}]
  \vspace{2pt}
  \begin{tabularx}{\linewidth}{X X}
    \textbf{Student Lead Signature:} & \textbf{Laboratory Instructor Approval:} \\[18pt]
    \rule{0.85\linewidth}{0.6pt} & \rule{0.85\linewidth}{0.6pt} \\
    \textbf{Elena Rostova} (ID: \#VIT-884219) & \textbf{Dr. Marcus Vance} (Lead Supervisor) \\
    Date: July 24, 2026 & Date: July 24, 2026
  \end{tabularx}
  \vspace{2pt}
\end{tcolorbox}

\end{document}
